% Part: first-order-logic
% Chapter: lindstrom
% Section: ls-property

\documentclass[../../../include/open-logic-section]{subfiles}

\begin{document}

\olfileid{mod}{lin}{lsp}

\olsection{સઘનતા અને લેવેનહાઇમ--સ્કોલેમ ગુણધર્મો}

હવે આપણે સઘનતા અને લેવેનહાઇમ--સ્કોલેમને અમૂર્ત તર્કશાસ્ત્રોના કિસ્સા
સુધીના સીધા વિસ્તારો આપીએ છીએ.

\begin{defn}
અમૂર્ત તર્કશાસ્ત્ર $\tuple{L, \models_L}$ પાસે \emph{સઘનતા
  ગુણધર્મ} છે, જો $L(\Lang{L})$-!!{sentence}sનો દરેક ગણ $\Gamma$ તેના
દરેક સાન્ત ઉપગણ $\Gamma_0 \subseteq \Gamma$ સંતોષ્ય હોય ત્યારે સંતોષ્ય
હોય.
\end{defn}

\begin{defn}
$\tuple{L, \models_L}$ પાસે \emph{અધોગામી લેવેનહાઇમ--સ્કોલેમ
  ગુણધર્મ} છે, જો દરેક સંતોષ્ય $\Gamma$ને !!a{enumerable} નિદર્શ હોય.
\end{defn}


\olref[bas][pis]{defn:partialisom}માંની આંશિક એકરૂપતાની વિભાવના તદ્દન
``બીજગણિતીય'' છે (એટલે કે, ભાષાનાં !!{sentence}sનો સંદર્ભ લીધા વિના,
માત્ર !!{structure}sની !!{language}~$\Lang{L}$ આપતા બિનતાર્કિક સંકેતો
વડે અપાયેલી છે), તેથી એ અમૂર્ત તર્કશાસ્ત્રોના કિસ્સામાં પણ લાગુ પડે છે.
પ્રથમ-ક્રમ તર્કશાસ્ત્રના કિસ્સામાં આપણે
\olref[bas][pis]{thm:p-isom2}થી જાણીએ છીએ કે બે !!{structure}s આંશિક
રીતે એકરૂપ હોય, તો તેઓ પ્રાથમિક રીતે સમકક્ષ છે. એ સાબિતી અમૂર્ત
તર્કશાસ્ત્રો સુધી ચાલતી નથી, કારણ કે મનસ્વી $!E \in L(\Lang{L})$ માટે
!!{formula}s પરનું આગમન ઉપલબ્ધ હોવું જરૂરી નથી. છતાં લેવેનહાઇમ--સ્કોલેમ
ગુણધર્મ હોય તો આ પ્રમેય સત્ય છે.
\sourcecorrection{OLLIN-004}{સ્થિર અંગ્રેજી મૂળની
  \texttt{ls-property.tex}, પંક્તિઓ~29--33માં આંશિક એકરૂપતા ભાષાનાં
  માત્ર અચળો પર આધાર રાખે છે એમ કહેવાયું છે. એની વ્યાખ્યામાં અચળો ઉપરાંત
  વિધેયો અને વિધેય-સંકેતોની શરતો પણ આવે છે; અહીં ``બિનતાર્કિક સંકેતો''
  કહીને યોગ્ય વ્યાપ રાખ્યો છે.}

\begin{thm}
\ollabel{thm:abstract-p-isom}
ધારો કે $\tuple{L, \models_L}$ એ લેવેનહાઇમ--સ્કોલેમ ગુણધર્મ ધરાવતું
સામાન્ય તર્કશાસ્ત્ર છે. ત્યારે કોઈ પણ બે આંશિક રીતે એકરૂપ
!!{structure}s $\tuple{L, \models_L}$માં પ્રાથમિક રીતે સમકક્ષ છે.
\end{thm}

\begin{proof}
ધારો કે $\Struct{M} \simeq_p \Struct{N}$, પણ કોઈ $!E$ માટે
$\Struct{M} \models_L !E$ જ્યારે $\Struct{N} \not\models_L !E$ પણ છે.
એકરૂપતા ગુણધર્મથી આપણે માની શકીએ છીએ કે $\Domain{M}$ અને $\Domain{N}$
અસંયુક્ત છે, અને વિસ્તાર ગુણધર્મથી માની શકીએ છીએ કે કોઈ સાન્ત
!!{language}~$\Lang{L}$ માટે $!E \in L(\Lang{L})$. $\Struct{M}$ અને
$\Struct{N}$ વચ્ચેની આંશિક એકરૂપતાઓનો ગણ $\mathcal{I}$ લો; અને કોઈ
સામાન્યતા ગુમાવ્યા વિના એમ પણ ધારો કે $p \in \mathcal{I}$ અને $q
\subseteq p$ હોય, તો $q \in \mathcal{I}$ પણ.

$\Domain{M}^{<\omega}$ એ $\Domain{M}$ના !!{element}sની બધી સાન્ત
શ્રેણીઓનો ગણ છે. $S$ને $\Domain{M}^{<\omega}$ પરનો જોડાણ દર્શાવતો
ત્રિસ્થાની સંબંધ ગણો; એટલે કે, $\mathbf{a}, \mathbf{b}, \mathbf{c}
\in \Domain{M}^{<\omega}$ હોય, તો $S(\mathbf{a}, \mathbf{b},
\mathbf{c})$ ત્યારે અને માત્ર ત્યારે જ્યારે $\mathbf{c}$ એ
$\mathbf{a}$ અને $\mathbf{b}$નું જોડાણ હોય. અને $T$ એવો ત્રિસ્થાની
સંબંધ ગણો કે $b \in M$ તથા $\mathbf{a}, \mathbf{c} \in
\Domain{M}^{<\omega}$ માટે $T(\mathbf{a}, b, \mathbf{c})$ ત્યારે અને
માત્ર ત્યારે જ્યારે $\mathbf{a} = a_1, \dots a_n$ અને
$\mathbf{c} = a_1, \dots a_n, b$. નવાં ત્રિસ્થાની !!{predicate}s
$P$ અને $Q$ લો અને એવી !!{structure} $\Struct{M}^*$ બનાવો જેનું
વ્યાપકક્ષેત્ર $\Domain{M} \cup \Domain{M}^{<\omega}$ છે, જેમાં
$\Struct{M}$ ઉપસંરચના છે અને જે $P$ તથા $Q$નું અર્થઘટન અનુક્રમે
જોડાણ-સંબંધો $S$ તથા $T$ તરીકે કરે છે (આથી $\Struct{M}^*$ !!{language}
$\Lang{L} \cup \{ P, Q\}$માં છે).

$\Domain{N}^{<\omega}$, $S'$, $T'$, $P'$, $Q'$ અને
$\Struct{N}^*$ને એ જ રીતે વ્યાખ્યાયિત કરો. પરિકલ્પનાથી $\Struct{M}
\simeq_p \Struct{N}$ હોવાથી $\Domain{M}^{<\omega}$ અને
$\Domain{N}^{<\omega}$ વચ્ચે એવો સંબંધ $I$ છે કે
$I(\mathbf{a}, \mathbf{b})$ ત્યારે અને માત્ર ત્યારે જ્યારે
$\mathbf{a}$ અને $\mathbf{b}$ એકરૂપ હોય અને
\olref[bas][pis]{defn:partialisom}નો આગળ-પાછળ ગુણધર્મ સંતોષે. હવે
$\Struct{A}$ને એવી !!{structure} ગણો જેનું !!{domain}
$\Struct{M}^*$ અને $\Struct{N}^*$નાં !!{domain}sનો સંઘ છે, જેમાં
$\Struct{M}^*$ અને $\Struct{N}^*$ ઉપ!!{structure}s છે, અને જેની
!!{language}માં એક વધારાનું દ્વિસ્થાની !!{predicate}~$R$ છે જેનું
અર્થઘટન સંબંધ~$I$ છે તથા !!{domain}s $\Domain{M}^*$ અને $\Domain{N}^*$ દર્શાવતાં
!!{predicate}s છે.
\sourcecorrection{OLLIN-005}{સ્થિર અંગ્રેજી મૂળની
  \texttt{ls-property.tex}, પંક્તિઓ~80--86માં નવી વ્યાપક સંરચનાને ફરી
  $\Struct M$ નામ અપાયું છે, જોકે એ જ અનુચ્છેદોમાં $\Struct M$ મૂળ
  આંતરિક સંરચના છે. નવી સંરચના માટે $\Struct A$ રાખીને બંને સંદર્ભો
  અલગ કર્યા છે; આ સુધારો આકૃતિ અને આગળની સાબિતીમાં પણ લાગુ પડે છે.}
\sourcecorrection{OLLIN-006}{સ્થિર અંગ્રેજી મૂળની
  \texttt{ls-property.tex}, પંક્તિ~86માં $\Domain{N}*$ લખવાથી તારક
  ઘાતમાં રહેતો નથી. અહીં સમાંતર $\Domain{M}^*$ પ્રમાણે
  $\Domain{N}^*$ લખ્યું છે.}

\begin{figure}[h]
  \centering
  \begin{tikzpicture}[node distance=2cm, auto, thick, >=stealth']
    \draw [rounded corners] (0,0) -- (8,0) -- (8,4) -- (0,4) --  cycle;
    \draw (2,2) circle (0.5cm);
    \draw (2,2) circle (1.25cm);
    \draw (6,2) circle (0.5cm);
    \draw (6,2) circle (1.25cm);
    \path node at (0.75,3.5) {\large $\Struct{A}$};
    \path node at (2,2) {\large $\Struct{M}$};
    \path node at (6,2) {\large $\Struct{N}$};
    \path node at (3.5,1) {\large $\Struct{M}^*$};
    \path node at (7.5,1) {\large $\Struct{N}^*$};
    \node (Idom) at (2.8,2) {};
    \node (Irng) at (5.2,2) {};
    \draw[<->, bend left] (Idom) to node {\large $I$} (Irng) ;
  \end{tikzpicture}
  \caption{આંતરિક આંશિક એકરૂપતા ધરાવતી !!{structure}~$\Struct{A}$.}
\end{figure}

નિર્ણાયક અવલોકન એ છે કે !!{structure}~$\Struct{A}$ની !!{language}માં
એવું $L$-!!{sentence} $!D_1$ છે જે $\Struct{A}$માં સાચું છે અને કહે છે
કે $\Struct{M} \models_L !E$ તથા $\Struct{N} \not\models_L !E$ (આમાં
સાપેક્ષીકરણ અને બૂલીય ગુણધર્મો જરૂરી છે); અને એવું પ્રથમ-ક્રમ
!!{sentence} $!D_2$ પણ છે જે $\Struct{A}$માં સાચું છે અને કહે છે કે
આંશિક એકરૂપતા~$I$ દ્વારા $\Struct{M} \simeq_p \Struct{N}$.
\sourcecorrection{OLLIN-007}{સ્થિર અંગ્રેજી મૂળની
  \texttt{ls-property.tex}, પંક્તિઓ~109--115માં $!D_1$ને પ્રથમ-ક્રમ
  વાક્ય કહેવામાં આવ્યું છે. મનસ્વી અમૂર્ત વાક્ય $!E$નો આંતરિક
  સંરચનાઓ પર સંતોષ સાપેક્ષીકરણથી $L$-વાક્ય તરીકે વ્યક્ત થાય છે;
  $!D_2$ પ્રથમ-ક્રમ વાક્ય છે અને સામાન્યતા પરથી $L$માં પણ આવે છે.
  અહીં આ પ્રકારભેદ સ્પષ્ટ કર્યો છે.}
લેવેનહાઇમ--સ્કોલેમ ગુણધર્મથી $!D_1$ અને $!D_2$ કોઈ
!!a{enumerable} નિદર્શ $\Struct{A}_0$માં એકસાથે સાચાં છે. તેમાં એવી
આંશિક રીતે એકરૂપ ઉપસંરચનાઓ $\Struct{M}_0$ અને $\Struct{N}_0$ છે કે
$\Struct{M}_0 \models_L !E$ અને $\Struct{N}_0 \not\models_L !E$.
\sourcecorrection{OLLIN-008}{સ્થિર અંગ્રેજી મૂળની
  \texttt{ls-property.tex}, પંક્તિઓ~115--118માં ગણનીય વ્યાપક નિદર્શ અને
  એની એક આંતરિક ઉપસંરચના બંનેને $\Struct M_0$ કહેવામાં આવ્યાં છે.
  વ્યાપક નિદર્શ માટે $\Struct A_0$ અને અંદરની બે સંરચનાઓ માટે
  $\Struct M_0,\Struct N_0$ રાખ્યાં છે.}
પરંતુ !!{enumerable} આંશિક રીતે એકરૂપ !!{structure}s વાસ્તવમાં
\olref[bas][pis]{thm:p-isom1}થી એકરૂપ છે, જે સામાન્ય અમૂર્ત
તર્કશાસ્ત્રોના એકરૂપતા ગુણધર્મનો વિરોધ કરે છે.
\end{proof}

\end{document}
